\chapter{Desarrollo}
\label{cap:desarrollo}

El capítulo anterior presenta los resultados que produce el framework
\framework. Este capítulo describe \emph{cómo se ha construido}: la
arquitectura software, las decisiones tecnológicas, la implementación
concreta del pipeline, la integración del modelo de lenguaje, la
interfaz web pública y las garantías de reproducibilidad. Es una lectura
complementaria a la Metodología del capítulo~\ref{cap:metodos}: aquélla
describía el diseño científico (qué se hace y por qué); ésta describe la
implementación de ingeniería (cómo se hace y con qué herramientas).

\section{Arquitectura general del framework}
\label{sec:dev_arquitectura}

\framework\ se organiza como una \textbf{tubería secuencial de nueve
pasos} implementados como módulos Python independientes que se ejecutan
desde un orquestador único. Cada paso lee sus entradas del sistema de
ficheros, produce artefactos intermedios (CSV, JSON, PDF de figuras) y
los deposita en rutas predecibles que el paso siguiente consume. Esta
disciplina de \textbf{artefactos persistentes} en disco tiene tres
consecuencias:

\begin{enumerate}
  \item Cualquier paso puede ejecutarse aislado, sin necesidad de
  recomputar los anteriores si sus salidas ya existen.
  \item Los resultados intermedios quedan disponibles para auditoría
  independiente sin necesidad de reejecutar el modelo.
  \item La interfaz web \framework\ es una capa de lectura sobre esos
  artefactos, lo que garantiza que \emph{lo que muestra la app coincide
  exactamente con lo que produjo el pipeline}, sin cifras hard-coded.
\end{enumerate}

La figura~\ref{fig:arquitectura} del anexo~\ref{anexo:interfaz} resume
esta separación entre pipeline (que produce artefactos) e interfaz (que
los consume).
Los pasos 1 y 2 realizan la ingesta técnica desde GEO; el paso 3 delega
la curación clínica al LLM; los pasos 4 a 6 hacen análisis diferencial
por cohorte y clasificación supervisada; los pasos 7 y 8 aplican
validación externa LODO y meta-análisis por consenso; el paso 9 produce
la firma final, el panel mínimo y las tablas de validación IHC. Un
décimo paso (paso 12 en la numeración interna del repositorio) despliega
la interfaz web sobre los artefactos producidos.

\section{Stack tecnológico}
\label{sec:dev_stack}

La totalidad del framework se implementa en \textbf{Python 3.12}, con
las librerías del ecosistema científico estándar como núcleo:

\begin{itemize}[leftmargin=1.2em]
  \item \texttt{pandas} y \texttt{numpy} para manipulación tabular y
  álgebra lineal.
  \item \texttt{scipy.stats} para las pruebas de hipótesis (Welch,
  Mann-Whitney) y el estadístico $d$ de Cohen.
  \item \texttt{statsmodels} para la corrección FDR de
  Benjamini-Hochberg y las regresiones adicionales.
  \item \texttt{scikit-learn} para los clasificadores supervisados
  (\texttt{LogisticRegression} con penalización L1, \texttt{RandomForestClassifier},
  \texttt{LinearSVC}), la validación cruzada anidada y la calibración
  con \texttt{CalibratedClassifierCV} (Platt e isotónica).
  \item \texttt{GEOparse} para el acceso programático a la API de GEO
  y a las tablas de anotación GPL empaquetadas con cada serie.
  \item \texttt{matplotlib} y \texttt{seaborn} para la generación de
  figuras estáticas.
\end{itemize}

Fuera del ecosistema de análisis, el framework se apoya en tres
componentes externos:

\begin{itemize}[leftmargin=1.2em]
  \item \textbf{Llama 3.3 de 70\,000 millones de parámetros} como
  modelo de lenguaje, servido mediante la API pública de \textbf{Groq}.
  \item \textbf{Streamlit} como framework de interfaz web
  multipágina para la aplicación \framework.
  \item \textbf{Git y GitHub} como sistema de control de versiones e
  infraestructura de publicación del código.
\end{itemize}

La decisión de usar exclusivamente software libre (con la única
excepción del acceso a la API de Groq, que es gratuita para el volumen
del proyecto y sustituible por un modelo local sin cambios estructurales)
persigue una reproducibilidad sin barreras económicas: cualquier
investigador con acceso a un ordenador estándar puede regenerar
íntegramente los resultados del TFM.

\section{Implementación del pipeline}
\label{sec:dev_pipeline}

Cada uno de los nueve pasos del pipeline es un script Python autónomo
en el directorio \texttt{agentes/}. Su nomenclatura sigue el patrón
\texttt{pasoN\_descripcion.py}, donde \texttt{N} preserva el orden
canónico de ejecución. A modo de referencia:

\begin{itemize}[leftmargin=1.2em]
  \item \texttt{paso1\_descarga.py} descarga cohortes GEO y aplica el
  mapeo sonda~$\to$~gen descrito en la sección~\ref{sec:descarga}.
  \item \texttt{paso2b\_bulk.py} normaliza cada cohorte (log$_2$ +
  cuantiles) y produce la matriz de expresión limpia.
  \item \texttt{paso3\_diferencial.py} orquesta la llamada al LLM para
  la curación de metadatos.
  \item \texttt{paso5\_graficos.py} y \texttt{paso6\_dashboard.py}
  ejecutan los análisis diferenciales y los agregan por cohorte.
  \item \texttt{paso7\_orquestador.py} y \texttt{paso13\_subtipo\_lodo.py}
  implementan el clasificador supervisado y su validación LODO.
  \item \texttt{paso20\_recalibracion.py},
  \texttt{paso21\_bootstrap\_ic.py}, \texttt{paso22\_comparativa\_ml.py}
  y \texttt{paso23\_validacion\_tcga.py} contienen los análisis
  adicionales de recalibración isotónica, intervalos de confianza
  bootstrap, comparativa entre modelos y validación en TCGA.
\end{itemize}

Todos los pasos aleatorizados fijan \texttt{random\_state=42} para
garantizar reproducibilidad estricta bit a bit. Los pasos leen sus
parámetros de un fichero de configuración YAML
(\texttt{configuracion/tareas.yaml}) que centraliza los umbrales
(por ejemplo, $|d| > 0{,}5$ para el consenso multi-cohorte), las rutas
de entrada y salida, y la lista canónica de cohortes participantes.

Adicionalmente, un paquete Python interno \texttt{tfm/} expone
utilidades comunes (carga de matrices, validación de estructuras
tabulares, cálculos de métricas) reutilizadas por varios pasos. Este
paquete se acompaña de una batería de \textbf{pruebas unitarias} en
\texttt{tests/} que verifica invariantes clave (dimensiones esperadas,
ausencia de valores nulos, alineamiento entre matriz de expresión y
metadatos).

\section{Integración del modelo de lenguaje}
\label{sec:dev_llm}

La curación clínica automática es la pieza más singular del framework y
merece un tratamiento propio. Se implementa en dos módulos:
\texttt{paso3\_diferencial.py} (orquestación) y \texttt{contexto\_tfm.py}
(prompts y control de calidad).

\subsection{Diseño del prompt}

El prompt sistema se construye como \textbf{tarea de clasificación de
texto acotada}, no como generación abierta. Sus componentes principales:

\begin{itemize}[leftmargin=1.2em]
  \item \textbf{Rol}: se define al modelo como ``anotador experto de
  datos biomédicos de expresión génica'', con especialización en cáncer
  de pulmón.
  \item \textbf{Instrucción explícita}: se le pide una etiqueta
  \emph{exactamente} de una lista cerrada (por ejemplo,
  \texttt{ADC}, \texttt{SQC}, \texttt{tumor\_indefinido}, \texttt{sano},
  \texttt{ambiguo}) según reglas concretas.
  \item \textbf{Contexto de muestra}: se pasan los campos
  \texttt{title}, \texttt{source\_name}, \texttt{characteristics\_ch1}
  y \texttt{description} del registro GEO original, sin
  preprocesamiento.
  \item \textbf{Regla de escape}: se instruye al modelo para responder
  \texttt{ambiguo} cuando el texto no permite una asignación segura, en
  lugar de forzar una etiqueta incorrecta. Esta regla es la que evita
  las alucinaciones y explica por qué la tasa de éxito no es del 100\,\%
  pero sí es \emph{honesta}.
\end{itemize}

\subsection{Control de coste y latencia}

El framework realiza $\sim$1\,400 llamadas al LLM (una por muestra
GEO). Sin agrupación, esto llevaría horas y consumiría cuota de forma
poco controlada. Groq permite una tasa de inferencia de decenas a
cientos de tokens por segundo por su hardware específico (LPU), lo que
hace viable la curación completa en pocos minutos y con coste marginal
gracias a su capa gratuita.

El módulo de orquestación implementa además:

\begin{itemize}[leftmargin=1.2em]
  \item \textbf{Caché en disco} de respuestas ya obtenidas, indexada
  por hash del prompt. Reejecutar el paso 3 no incurre coste si los
  metadatos de entrada no cambian.
  \item \textbf{Reintentos con backoff exponencial} ante fallos
  transitorios de la API.
  \item \textbf{Auditoría manual} de un subconjunto (38 muestras
  seleccionadas aleatoriamente y etiquetadas manualmente por el autor)
  usada por \texttt{paso24\_auditoria\_manual\_llm.py} para estimar la
  precisión real de la curación.
\end{itemize}

\subsection{Aislamiento del LLM del análisis estadístico}

Una decisión de diseño explícita: el LLM \emph{sólo} etiqueta muestras.
No participa en la selección de genes, ni en la clasificación, ni en la
evaluación de rendimiento. Todos los pasos posteriores son
deterministas dado el conjunto de etiquetas producido por el LLM. Esta
separación evita que sesgos paramétricos del modelo se propaguen al
análisis y facilita auditoría independiente.

\section{Interfaz web \framework}
\label{sec:dev_interfaz}

La aplicación web \framework\ está construida en Streamlit y se
estructura como un dashboard multipágina que expone el trabajo a
audiencias tanto técnicas como divulgativas. El código de la app reside
en \texttt{agentes/paso12\_web\_chatbot.py} (router principal) y
\texttt{agentes/paso12\_dashboard.py} (dashboard analítico).

\subsection{Arquitectura multipágina}

La aplicación se organiza en siete capítulos navegables:

\begin{enumerate}[leftmargin=1.4em]
  \item Portada y presentación del framework.
  \item Introducción y objetivos.
  \item Metodología resumida.
  \item Resultados globales con buscador de gen en tiempo real.
  \item Vista detallada por cohorte (PCA, volcano plot, heatmap).
  \item Conclusiones y biomarcadores por linaje.
  \item Asistente conversacional Llama 3.3-70b acotado a los
  resultados del trabajo.
\end{enumerate}

Cada página \textbf{lee los datos en tiempo de renderizado} desde los
CSV/JSON producidos por el pipeline. Ningún valor está hard-coded en el
código de la app: si el pipeline se reejecuta con nuevas cohortes o
parámetros, la app refleja los nuevos números automáticamente al
recargar la página.

\subsection{Asistente conversacional acotado}

La página del asistente merece mención por su diseño defensivo contra
alucinaciones. El módulo \texttt{contexto\_tfm.py} construye
dinámicamente un prompt sistema que:

\begin{itemize}[leftmargin=1.2em]
  \item Incluye los ficheros de resultados como contexto (firma,
  panel, métricas LODO, tabla IHC, validación TCGA).
  \item Fija instrucciones estrictas: el modelo debe responder
  \emph{únicamente} con información derivable de los ficheros
  proporcionados. Si una pregunta cae fuera del contexto, debe
  declararlo explícitamente.
  \item Distingue explícitamente entre genes del panel mínimo, genes
  de la firma amplia y genes no incluidos.
\end{itemize}

El resultado es un asistente que puede responder preguntas del tipo
``\emph{¿qué biomarcadores identifica el framework para carcinoma
escamoso?}'' con las cifras y genes correctos del trabajo, pero que
declina responder preguntas de conocimiento general sobre oncología
que no estén respaldadas por los datos.

\section{Reproducibilidad, tests y despliegue}
\label{sec:dev_reproducibilidad}

\subsection{Garantías de reproducibilidad}

Además del \texttt{random\_state=42} y de los artefactos versionados,
el framework declara sus dependencias exactas en
\texttt{requirements.txt} con versiones fijadas. El proyecto usa
\texttt{pyproject.toml} para definir el paquete \texttt{tfm/} como
instalable, lo que garantiza que las importaciones dentro del paquete
funcionan en cualquier entorno donde se ejecute
\texttt{pip install -e .}.

\subsection{Pruebas automatizadas}

El directorio \texttt{tests/} contiene pruebas unitarias con
\texttt{pytest} que verifican:

\begin{itemize}[leftmargin=1.2em]
  \item Alineamiento entre matriz de expresión y metadatos (sin
  muestras huérfanas).
  \item Consistencia de dimensiones tras cada paso del pipeline.
  \item Cumplimiento de los criterios de inclusión de cohortes.
  \item Determinismo del clasificador entre ejecuciones con la misma
  semilla.
\end{itemize}

Estas pruebas se ejecutan localmente y sirven como red de seguridad
frente a regresiones al modificar el código.

\subsection{Repositorio público y despliegue}

El código fuente completo se publica en GitHub bajo el repositorio
\url{https://github.com/danitapiadiez-gif/TFM-Bioinformatica}. La
organización de directorios sigue la convención estándar:

\begin{itemize}[leftmargin=1.2em]
  \item \texttt{agentes/} --- los nueve pasos del pipeline y la app web.
  \item \texttt{tfm/} --- paquete de utilidades reutilizables.
  \item \texttt{tests/} --- suite de pruebas automatizadas.
  \item \texttt{configuracion/} --- parámetros centralizados en YAML.
  \item \texttt{memoria/} --- fuente LaTeX de esta memoria.
  \item \texttt{presentacion/} --- fuente de la defensa.
\end{itemize}

La aplicación web puede desplegarse en \textbf{Streamlit Community
Cloud} conectando el repositorio directamente. El único secreto externo
requerido es la clave de API de Groq, que se inyecta como variable de
entorno y nunca se versiona. Esta configuración permite que el tribunal
o cualquier tercero interesado pueda desplegar su propia instancia de
la aplicación con un esfuerzo mínimo.

Los anexos~\ref{anexo:repo} (estructura del repositorio),
\ref{anexo:interfaz} (interfaz web) y~\ref{anexo:despliegue} (guía de
despliegue) documentan exhaustivamente los aspectos técnicos que en
este capítulo se han presentado de forma sintética.
